\section{Related Work}
\label{sec:relatedworks}

\paragraph{\Ematching}

The \ematching problem is first described by a seminal thesis of
 \citet{nelson-thesis} in the context of program verification,
 although it has been proved earlier that the problem
 of \ematching is NP-complete.
\citet{simplify} shows the first formal description of
 an \ematching algorithm, which is implemented as
 backtracking iterators.
Modern SMT solvers implement various \ematching algorithms
 for quantifier instantiation.
Moreover, equality saturation-based optimizers uses \ematching to discover
 new terms equivalent to the original program.
Several improvements are introduced
 later~\cite{moskal, efficient-ematching},
 all based on the idea of backtracking.
In particular, \citet{efficient-ematching} describes an
 efficient \ematching virtual machine that computes
 \ematching incrementally.
Their algorithm uses several indices to track changes
 incrementally and speed up \ematching.
However, their algorithm is difficult to implement and
 cannot utilize the equality constraints in general.
Their incremental maintenance algorithm can also be viewed
 as a specialized algorithm for incremental view maintenance
 if taking a relational view of \egraphs.

\paragraph{Graph processing engine}

\Ematching bears many similarity with graph analytics queries
 like graph pattern matching, both of which try to find patterns
 on an \egraph.
Since many real-world data such as social networks
 can be naturally expressed as graphs, systems for graph analytics
 has increasing becoming the focus of database researches
 and many graph processing engines are developed.
Some engines use dedicated graph data structures and
 graph algorithms for graph analytics~\cite{ligra, snap, powergraph}.
Other engines use a relational database system
 for performing the analytics~\cite{socialite,grail,emptyheaded}.
In particular, Emptyheaded shows that a relational engine using
 the generic join algorithm can beat low-level systems dedicated
 to graph analytics~\cite{emptyheaded}.
Consistent with this conclusion, we use the generic join algorithm for
 the \ematching problem.


\paragraph{Equality saturation}

\paragraph{SMT solver}

\paragraph{Rewrite systems}

Rete's algorithm.

\paragraph{Join algorithms}

%%% Local Variables:
%%% mode: latex
%%% TeX-master: "main"
%%% End:
